\documentclass[11pt,spanish]{article}
%\usepackage[latin1]{inputenc}
\usepackage[utf8]{inputenc}
\usepackage[spanish,es-nodecimaldot]{babel}
\usepackage{graphicx}
\usepackage{tikz}
\usepackage{amsmath,amssymb}


\oddsidemargin=0.5cm
\textwidth=16cm
\textheight=48\baselineskip
\topmargin=-1.5cm
\parskip 2truemm
\parindent 0pt
\newcounter{problema}
\typeout{****************************************}
\typeout{****** Configuracion de Guia 11pt ******}
\typeout{****************************************}
%%%%%%%%%%%%%%%%%%%%%%%%%%%%%%%%%%%%%%%%%%%%%%%%%%%%%%%%%%%%%%%%%%%%%%%

\begin{document}

\renewcommand{\labelenumi}{\em \alph{enumi})}
\renewcommand{\labelenumii}{\em \arabic{enumii})}
\newcommand\al{\alpha}
\newcommand\n{\nabla}
\newcommand\F{\vec F}
\newcommand\G{\vec G}
\newcommand\noi{\noindent}
\newcommand\re{\mbox{Re}}
\newcommand\im{\mbox{Im}}
%\newcommand\arg{\mbox{arg}}
\newcommand\Arg{\mbox{Arg}}
\newcommand{\prob}{\stepcounter{problema} \noindent{\bf Problema
\arabic{problema}: }}
\newcommand\pin[2]{\big\langle #1,#2 \big\rangle}
\newcommand\pd[2]{\frac{\partial #1}{\partial #2}}
\newcommand\R{\mathbb{R}}

%\input{/home/ortiz/lib/definiciones.tex}
\begin{center}
{\large\sc Facultad de Matem\'atica, Astronom\'ia, F\'isica y Computaci\'on, U.N.C.}

\vspace{2mm}
{\large M\'etodos Matem\'aticos de la F\'isica II}

\vspace{2mm}
{\large Gu\'ia N$^{\sf o}$ 6 (2021)}
\end{center}

\prob{} Verificar que 
$$
u_p(x,t) = \frac{1}{2c} \int_0^t\int_{x-c(t-\tau)}^{x+c(t-\tau)}
\psi(s,\tau)~ds~d\tau.
$$
satisface el problema
$$
\pd{^2u}{t^2} - c^2\pd{^2u}{x^2} = \psi(x,t), \qquad u:D\subset\R^2\to
\R,
$$
con dato inicial cero.




\prob{}
Resolver la ecuaci\'on $u_{tt}=c^2 u_{xx}$ con las siguientes condiciones:

\begin{enumerate}
\item $u(x,0)=0$ y $u_t(x,0)=\sin x$

\item $u(x,0)=e^x$ y $u_t(x,0)=\cos x$

\end{enumerate}


\prob{} Probar que si los datos iniciales $f(x)$ y $g(x)$ tienen soporte
compacto, o decaen junto con sus derivadas primeras suficientemente rápido a
cero cuando $|x|\to\infty$, y el problema 
\begin{equation*}\begin{split}
\pd{^2u}{t^2} - c^2\pd{^2u}{x^2} =& 0, \qquad u:D\subset\R^2\to
\R,\\
u(x,0) =& f(x),\\
\pd{u}{t}(x,0) =& g(x).
\end{split}
\end{equation*}
conserva la energía de la solución. {\em Sugerencia:} Muestre que la derivada de
$E(t)$ respecto al tiempo es cero, y por lo tanto $E(t) = E(0)$ para todo $t>0.$

\prob{} Resolver el problema de las oscilaciones transversales propias de una cuerda homogenea de longitud $L$ si:

\begin{enumerate}
\item Los extremos de la cuerda estan fijos r\'igidamente. 

\item los extremos de la cuerda estan libres. Es decir, $\cfrac{du}{dx}=0$ en los extremos de la cuerda. Esto tiene lugas cuando los extremos de la cuerda est\'an sujetos mediante anillos (de masas despreciables), que se deslizan sin rozamiento sobre barras paralelas.


\end{enumerate}
\prob{} Defina la Energía para el problema 
\begin{equation*}\begin{split}
&\pd{^2u}{t^2} - c^2\pd{^2u}{x^2} = 0, \qquad u:[0,L]\times[0,\infty)\subset\R^2\to
\R,\\
&u(0,x) = f(x),\quad  \pd{u}{t}(0,x) = g(x),\qquad \mbox{(datos iniciales)}\\
&u(0,t) = u(L,t)=  0, \qquad t\ge 0, \qquad \mbox{(condiciones de contorno)}
\end{split}
\end{equation*}
y muestre que esta es constante en el tiempo.

\prob{} Condidere el problema para la cuerda de guitarra, igual al anterior
salvo por la introducción de una pequeña disipación, representada mediante un término de derivada temporal primera. La ecuación diferencial es entonces
\[
\pd{^2u}{t^2} - c^2\pd{^2u}{x^2} + \gamma \pd{u}{t} = 0, \qquad 0 < \gamma <
\frac{2\pi c}{L}.
\]
Resuelva el problema. Muestre que la solución y su energía en este caso decaen
con el tiempo. >Con qué ley temporal?

\prob{} Hallar a solución de la ecuación de ondas homogénea en la franja
$(x,t) \in D = [0,1]\times [0,\infty),$ con condiciones iniciales homogéneas, es
decir
\[
u(x,0) = 0, \qquad \pd{u}{t}(x,0) = 0,
\]
y con condiciones de contorno 

\begin{equation}\label{bc_frontera_izq}
\pd{u}{t}(x_0,t) + c\pd{u}{x}(x_0,t) = h_0(t).
\end{equation}
en la frontera $x=0$ con
\[
h_0(t) = \begin{cases}
1 &  0<t\le 1,\\
0 &  1<t,
\end{cases}
\]
y condición de reflección 
\[
\pd{u}{x}(1,t) = 0
\]
sobre la frontera $x=1.$

\prob{} Resuelva la ecuaci\'on de ondas, $c^2 \Delta u = u_{tt}$, para una membrana circular de radio $R$ con el borde fijo sabiendo que $u(r,\theta,0)=0$ y $u_t(r,\theta,0)=r \sen (3\theta)$. 

\prob{} Resolver la ecuaci\'on de onda de una membrana en forma de cu\~na que se encuentra inicialmente en reposo con las condiciones de borde tal como muestra la figura.  En $t=0$ la membrana tiene la forma $f(r,\theta,0)=r\sin(6 \theta)$.
\begin{center}
\begin{tikzpicture}[scale=1.2]
  \filldraw[fill=gray!40!white, draw=black!50!black]
    (2,0) -- (4,0) arc (0:35:4) -- (0,0);
\draw[<-] (5,0) -- (0,0); 
\draw[<-] (0,2.5) -- (0,0);
\node at (5,-0.2) {\it \large x};
\node at (-0.2,2.5) {\it \large y};
\node at (4,-0.2) {\it \large a};
\node at (2.7,0.8) {$c^2 \Delta u =u_{tt}$};
\node at (4.3,1.2) {$u=0$};
\node [rotate=35 ]  at (1.5,1.3) {$u=0$};
\node at (2,-0.2) {$u=0$};
\draw [<-] (1,0) arc (0:60:16.5pt);
\node at (1.4,0.25) {$\alpha=\frac{\pi}{3}$};
\end{tikzpicture}   
\end{center}


\prob{} Considere la  ecuaci\'on de ondas $u_{tt}= c^2\Delta u$ en el interior de una esfera de radio $a$ tal que $u=u(r,t)$. 

a) Discuta una transformaci\'on de tipo $u = r^{\alpha}\psi$ para reducir esta ecuaci\'on a la  ecuaci\'on del onda usual.

b) Resuelva la ecuación anterior que satisface las conidicones $u(r=a,t)=0$, $u(r,0)=\beta$ y $u_t(r,0)=0$, donde $\beta$ es una constante.

\prob{} Cierta cuerda de largo $L$ cuya densidad es variable est\'a tensada por sus extremos que est\'an fijos y sus desplazamientos $u$ (respecto de la posici\'on de equilibrio) est\'an bien modelados por la ecuaci\'on (de onda)
\[ u_{tt} = c^2 (1+ \epsilon x )^2 u_{xx}\;,\;\; 0<x<L\;,\]
donde $\epsilon >0$.
Determine los modos de esta cuerda.\\
\textit{Sugerencia:} Use la transformaci\'on de variables  $y:= 1+ \epsilon x$ para reescribir la ecuaci\'on.\\



\end{document}



